\section{Synthesis and Research Agenda}
\label{sec:agenda}

\subsection{What the literature currently supports}

The mapped evidence supports six qualified conclusions.

First, tactile-specific pretraining is useful. Across robotic play,
self-supervision, multimodal alignment, and cross-sensor learning, tactile
representations outperform purely task-specific training in several downstream
settings~\cite{guzey2023tdex,higuera2025sparsh,yang2024unitouch,feng2025anytouch}.
This justifies reusing a strong local encoder rather than learning raw tactile
features from scratch for every bimanual experiment.

Second, spatial and relational inductive biases can be useful when the task has
known structure. Morphology graphs, tactile graphs, object-centric dynamics, and
contact fields have succeeded in in-hand manipulation, physics prediction,
packing, and contact tracking~\cite{funabashi2022distributed,yang2023tacgnn,ai2024robopack,higuera2023ncf}.
This establishes plausibility, not universal superiority over attention.

Third, bimanual dexterity has become an active, increasingly crowded research
area. Large simulation suites, demonstration-generation systems, retargeting, and
real policies now cover many complex tasks
~\cite{chen2022bidexhands,shaw2024bidex,jiang2024dexmimicgen,li2025maniptrans}.
Novelty must therefore be located in a specific mechanism or evaluation gap.

Fourth, contact-centric and tactile-reactive control is moving toward broader
embodiments, tools, and data scales~\cite{niu2026trex,yuan2026ftp1,ma2026semanticcontact,madan2026tactic,zhou2026touchworld}.
Because many examples are recent preprints or programme-listed works, they define
a competitive frontier more than a settled empirical consensus.

Fifth, the strongest competing explanations are increasingly structured even
when they are not tactile graphs. Directed tactile prediction
~\cite{xue2026hitacwam}, mesh-prior plus diffusion-residual dynamics
~\cite{wang2026meshpriordit}, temporal-flow supervision
~\cite{yang2026temporalflowvla}, prediction-error recurrence
~\cite{sawada2026predvla}, and failure-aware recovery
~\cite{zhao2026flare} each isolate a different source of performance. This makes
matched objective, history, and recovery controls part of the minimum evidence
for an explicit contact-structure claim.

Sixth, direct evidence now reaches beyond isolated local tactile features.
BiView-Touch supports useful cross-hand predictive dependence in human recordings;
causal force-history plus compliance supports real bimanual release/hold decisions;
and CADeT supports visually inferred transmission modes with active probing
~\cite{liang2026biviewtouch,marra2026temporaltactilehandover,hu2026cadet}.
These results narrow the gap. They do not jointly demonstrate independent robot
load-path validity, matched structural benefit, calibrated uncertainty, and
held-out contact-participant topology. Meanwhile, ContactWorld and the October
SimpleTouch/VISTA controls make representation, goal information, and feedback
frequency separate competing explanations
~\cite{zhang2026contactworld,yang2026simpletouch,wong2026vista}.

\subsection{What the literature does not yet establish}

The evidence does not yet establish that an inferred dynamic contact graph is
better than a matched temporal Transformer for bimanual load transfer. It does not
establish that local tactile foundation models preserve cross-interface wrench
information. It does not yet establish robust, uncertainty-tested recovery of robotic
load-path topology from deployable touch and proprioception across changed
participant structures. It also does not provide a
standard benchmark for transfer from handover to cooperative tool use.

These gaps are linked by an evidence chain:

\begin{equation}
\begin{aligned}
\text{sensor signal} &\rightarrow \text{local contact factor}
\rightarrow \text{entity/interface identity}\\
&\rightarrow \text{dynamic relation}
\rightarrow \text{future physical outcome}
\rightarrow \text{control decision}.
\end{aligned}
\label{eq:evidence_chain}
\end{equation}

Current papers often validate one or two adjacent links. A local encoder validates
the first link; a privileged graph policy validates the last links conditional on
correct structure; a successful bimanual imitation policy validates behaviour but
may leave the intermediate mechanism unidentified. A high-impact study need not
solve the complete chain, but it should state exactly which link it tests.

\subsection{Eight priority problems}

\begin{longtable}{P{0.06\linewidth} P{0.25\linewidth} P{0.30\linewidth} P{0.29\linewidth}}
\caption{Prioritized research agenda for structured bimanual tactile
representations.}\label{tab:agenda}\\
\toprule
\textbf{P} & \textbf{Problem} & \textbf{Decisive experiment} & \textbf{Potential deliverable} \\
\midrule
\endfirsthead
\toprule
\textbf{P} & \textbf{Problem} & \textbf{Decisive experiment} & \textbf{Potential deliverable} \\
\midrule
\endhead
1 & Observation-grounded dynamic interfaces & infer hand--object edges from tactile, proprioception, and optional vision; compare with oracle and random edges & contact-inference benchmark and calibrated edge model \\
2 & Cross-interface physical targets & predict phase, load share, future wrench, slip, and relaxation state across short horizons and sustained contact & probe suite and physically annotated handover data \\
3 & Matched structural comparison & compare concat, attention, directed/flat prediction, prediction-error recurrence, graph-only, residual-only, inferred, and oracle structure with equal budgets & defensible evidence for or against explicit structure \\
4 & Cross-topology transfer & handover pretraining followed by tool--workpiece adaptation under label/data budgets & compositional transfer method and benchmark \\
5 & Uncertainty under sensor failure & missing taxels, saturation, wear, thermal/viscoelastic drift, latency, and sensor-type shift & risk-aware contact representation \\
6 & Privileged-to-observation distillation & train with simulator contacts/wrenches; deploy an observation-only student & scalable supervision without deployment leakage \\
7 & Causal structure tests & role swap, receiver delay, history-order/edge rewiring, physical perturbation, and retry/reset separation & evidence that relations mediate predictions/control and recovery \\
8 & Interoperable data lineage & common schema, split manifests, timestamps, frames, masks, release versions, and label provenance & reproducible BiLoad-style evidence matrix and data tooling \\
\bottomrule
\end{longtable}

\subsection{A staged experimental programme}

The literature suggests a risk-reducing order of work.

\paragraph{Stage 1: audit and offline probing.}
If a downloadable, licensed T-Rex subset is verified, audit 20--100 episodes for
timestamps, missing channels, task phases, object identity, and label provenance
~\cite{niu2026trex}. Otherwise, make the VTDexManip bimanual handover task the
first executable benchmark while T-Rex ingestion remains a parallel artifact
audit~\cite{liu2025vtdexmanip}. Reproduce PhysGraph as stock public code and a separately
named corrected-bias variant; do not merge the two interpretations
~\cite{li2026physgraph}. The output should be a unified schema and an initial
proprio-only/concat/temporal probing table, not a large control policy.

\paragraph{Stage 2: representation isolation.}
Fix one local tactile encoder and compare matched fusion/structure variants on
phase, future edge, load share, future interface wrench, slip risk, and
uncertainty. The matrix should include a flat versus directed tactile hierarchy,
prediction-error recurrence, training-only temporal grounding, graph-only versus
global-residual-only models, and their structured combination
~\cite{xue2026hitacwam,wang2026meshpriordit,yang2026temporalflowvla,sawada2026predvla}.
Match observation refresh and executed action horizon as well
as representation budgets. Include independent load labels and a dense cross-hand
fusion control; neither phase proxies nor attention maps establish force transmission.
Use episode-, object-, and contact-sequence splits, plus contact-duration and
temperature strata where sustained loading is present. Early project
decisions may use three seeds, but final claims should use five or more where
variance warrants. Internal go/no-go thresholds may guide resource allocation,
yet they must not be written as achieved results or retrospectively optimized
paper hypotheses.

\paragraph{Stage 3: closed-loop handover.}
Only after the offline model beats strong temporal baselines should it enter
closed-loop quasi-static handover. Compare fixed schedules, causal force-history with compliance, unstructured
tactile feedback, anchored tokens, cross-hand completion, and inferred structure
~\cite{marra2026temporaltactilehandover,liang2026biviewtouch}. Evaluate safe release, peak
force, drop rate, time, uncertainty, failure detection, local retry, and reset
under receiver delay and object-property changes. Recovery branches should be
reported separately so that controller recovery does not masquerade as a
representation improvement~\cite{zhao2026flare}.

\paragraph{Stage 4: gated cooperative tool use.}
Freeze a rigid stylus or wiping-puck--plate task early, but move to full closed-loop
experiments only if (i) representation gains survive matched budgets and OOD
splits, and (ii) the coordination-necessity check confirms that the stabilizing
hand materially affects the outcome. The second study should emphasize transfer
from a three-node handover path to a longer tool--workpiece path. If the gate
fails, a valuable pivot is privileged-to-observation graph distillation or an
interface-inference benchmark rather than an underdiagnosed hardware demo.

\subsection{Where a bimanual contact-graph encoder can contribute}

An encoder built around bimanual contact flow remains scientifically meaningful,
but only under a narrow and testable formulation. Its strong contribution would
not be ``a GNN for two tactile hands.'' It would be the combination of:

\begin{itemize}
  \item hierarchical sensor/link/entity nodes that preserve quality masks and
  coordinate frames;
  \item dynamic, probabilistic edges inferred from deployable observation;
  \item object- or task-frame messages that represent load propagation through
  objects and tools rather than a permanent left--right shortcut;
  \item self-supervised or privileged-distilled edge/physics targets;
  \item calibrated uncertainty and explicit missing-sensor behaviour; and
  \item matched-budget evidence that the factorization improves transitions,
  data efficiency, or cross-topology generalization.
\end{itemize}

Under this formulation, the encoder is a major hypothesis and potential
contribution. Without observation-grounded edges and decisive comparisons, it is
an implementation variant in a crowded graph-and-touch landscape.

\subsection{Publication opportunities}

The research agenda naturally separates into two publishable units. The first is
a benchmark-and-representation paper: synchronized handover data and probes,
matched structural baselines, dynamic interface inference, and careful
observability audits. The second is a transfer paper: reuse of interface/relation
functions from handover in cooperative tool use, with role swap, new nodes,
limited-data adaptation, and causal perturbations. A third, lower-hardware-risk
alternative is an artifact paper on privileged-to-observation contact inference
and reproducible code auditing across simulators.

The literature review itself could support publication after a more systematic
update: preregistered database searches, deduplication across indexes, dual
screening, complete author/publication metadata, code/data-availability coding,
and a public evidence table. The current manuscript supplies the conceptual
taxonomy and reproducible local scaffold, but it should not be labelled a
systematic review until those steps are completed.
